\section{引言}
\label{sec:introduction}

知识图谱（KG）用三元组保存事实。每条三元组包含主体、关系和值。例如，收据图谱可以保存收据编号、字段“total”和取值“120.00”。从文本构建的图谱常有值缺失、值错误、三元组重复或关系不合法等问题\citep{ref_ji2021,ref_paulheim2017}。修复需要先找到这些错误，再选择与原文相符的修改。

第一个问题是决定检查什么。三元组重复、字段不合法和值错误的原因各不相同。质量检查工具覆盖结构、规则和事实内容\citep{ref_zaveri2016,ref_kontokostas2014,ref_wang2021}。我们按实体、图和上下文三个范围组织检查。它们查找重复事实，检查文档编号和字段名，并将取值与原文匹配。这些结果共同组成评估概况，记录哪些位置需要修复。

第二个问题是决定保留哪些修改建议。模型可以阅读文本并提出取值，但每个字段也有相应规则。SHACL 用于表达图规则，近期修复方法将这些规则加入大语言模型提示\citep{ref_shacl2017,ref_lin2025}。我们用来源和字段约束检查生成的三元组。过滤器在允许的关系和字段数量范围内，保留与来源匹配的值，并记录每条被拒绝三元组的原因。

第三个问题是在修复动作中作出选择。多个建议可能针对同一字段，一次编辑也可能影响其他检查。规则、学习模型和源文本提供了不同的图修复方式\citep{ref_paulheim2017,ref_pan2023}。我们测试尺度感知路由，用评估概况为各类编辑赋予权重。选择器在图副本上试改候选，检查边界，并为每项编辑打分。它执行得分最高且为正的允许编辑，然后再次检查图。

本文任务使用保存单篇文档字段的图。每张图具有固定文档编号和允许关系名单。每种关系最多对应一个值。主修复路径先清理输入，再根据原文生成完整图，最后过滤其中的三元组。错误报告或来源行提示可以引导模型。我们分别为这条路径和顺序选择器设置对照，并保存各自的执行日志。

实验使用三个中文记录集合、SROIE 收据和 CUAD 合同。我们用 20 张新收据设置来源行索引，然后冻结设计，在另外 60 张收据上测试。合同测试使用 56 份文档。两名评审检查 200 项参考差异和 200 项系统编辑。共享提示以及同一输出过滤前后的评分展示各步骤的作用。实验衡量字段恢复、事实保留、请求成本和人工编辑接受度。

本文有三项贡献：
\begin{itemize}
    \item 将局部、图和来源检查汇入评估概况，在字段约束下引导基于来源的修复。
    \item 设计输出过滤流程，并单独测试尺度感知选择器，比较试改收益、成本和图边界。
    \item 在注入缺陷、模型输出、收据和合同上比较组件，开展人工评价，并公开输出和决策日志。
\end{itemize}

第~\ref{sec:related_work}~节回顾相关工作。第~\ref{sec:framework}、\ref{sec:implementation}~节介绍方法与实现。第~\ref{sec:experiments}~节报告实验和结果。第~\ref{sec:conclusion}~节总结全文。